% !PRINTED 678
\weberpage{678}
%<<CONT
and continue this sequence as long as one can still find functions $\lambda_s$ that are not contained in the form (3).

By 4., 3) the functions $\lambda_1, \lambda_2, \ldots, \lambda_s$ are linearly independent in the sense that no linear equation with rational coefficients in $z$ subsists between them. Since there are not more than $n$ linearly independent functions, $s$ cannot be greater than $n$. On the other hand, as long as $s \leqq n$, there are always still integral functions that are not linearly dependent on $\lambda_1, \lambda_2, \ldots, \lambda_{s-1}$, and among them also ones of least exponent. We therefore obtain $n$ functions
\weq{4}{L = (\lambda_1, \lambda_2, \ldots, \lambda_n),}
and these form a basis of the field $\Omega$ with respect to $z$. They form a minimal basis; for if this were not the case, then for some $s \leqq n$ a function $x_1\lambda_1 + x_2\lambda_2 + \cdots + x_s\lambda_s$ would have to be divisible by a linear function $z - c$ without $x_s$ being divisible by $z - c$. If one reduces $x_1, x_2, \ldots, x_s$ to their residues with respect to $z - c$, an equation results:
\[
c_1\lambda_1 + c_2\lambda_2 + \cdots + c_s\lambda_s = (z - c)\mu,
\]
in which the $c_1, c_2, \ldots, c_s$ are constant, $c_s$ different from zero, and $\mu$ an integral function. Now $\mu$ is not contained in the form (3), and by 2. and 3. its exponent is smaller than the exponent of $\lambda_s$, which contradicts the determination 4. 3). Hence:

5. \gesperrt{The functions $\lambda_1, \lambda_2, \ldots, \lambda_n$ form a minimal basis of $\Omega$ with respect to $z$. It is called a normal basis.}

The exponents $r_1, r_2, r_3, \ldots, r_n$ of these functions satisfy the condition:
\weq{5}{r_1 = 0, \quad 1 \leqq r_2 \leqq r_3 \ldots \leqq r_n.}
For if $r_s < r_{s-1}$, then by 3. 3) one would have had to take $\lambda_s$ in place of $\lambda_{s-1}$.
% [sic] printed "3. 3)", the rule invoked (taking the function of least exponent) is the recurrent determination 4. 3), cited as such twice earlier on this page

Of the functions $\lambda_i$ none is divisible by $z$; for if $\lambda_s = z\mu$, then $\mu$ would have a smaller exponent than $\lambda_s$ and would not be contained in the form (3). Thus $\mu$ would have to take the place of $\lambda_s$.

6. Accordingly (by 1.)
\[
r_1, r_2, \ldots, r_n
\]
are the exponents of
\weq{6}{\lambda_1' = z'^{r_1}\lambda_1, \quad \lambda_2' = z'^{r_2}\lambda_2, \ldots, \lambda_n' = z'^{r_n}\lambda_n}
in regard to $z' = 1:z$, and these functions are the elements of a normal basis of $\Omega$ with respect to $z'$.
